%% =====================================================================
%%  T-15-04  The Fantasy Floor
%%  -------------------------------------------------------------------
%%  One idea per file. Core is mandatory. Slots in canonical order.
%%  Max 700 words. No layout commands. No filler. New source material
%%  for this entry goes in sources/B3/T-15-04.tex -- never by
%%  rewriting this file (docs/RULES.md R0.1).
%%  =====================================================================
%% @kind: topic
%% @id: T-15-04
%% @title: The Fantasy Floor
%% @subtitle: The truth is ugly --- whoever manufactures pleasantness is an oasis, and priced like one
%% @chapter: c15
%% @order: 40
%% @status: stable
%% @sources: B3
%% @updated: 2026-09-19

\topic{T-15-04}{The Fantasy Floor}{The truth is ugly --- whoever manufactures pleasantness is an oasis, and priced like one}

\begin{core}
People avoid the truth when it is ugly, and much of the truth is ugly. Offer no truths you are not prepared to be hated for --- disenchantment arrives with anger attached: the person who manufactures romance or conjures fantasy is an oasis in the desert, and everyone flocks. This entry states the appeal plainly --- and states, equally plainly, the bill the oasis pays.
\end{core}
\begin{mechanism}
B3's law: life is so harsh and distressing that those who can manufacture romance are like water in a dry country; there is great power in tapping the fantasies of the masses. The mechanism is demand arithmetic, not deception: reality under-delivers against the images people carry, so a supplier of better-feeling material takes market share from suppliers of accurate material --- and cannot be competed against on accuracy, which is not the currency. The fantasy supplier sells absence-of-pain at premium: the alchemist selling transformation, the consultant selling the inevitable triumph, the lover selling the exceptional. Two failure triggers are built in and must be stated: first, disenchantment rage --- the audience that discovers the oasis was staged does not merely leave, it attacks, because the fantasy was purchased with hope and hope's refund is demanded in malice; second, the source's closing concession --- fantasy need not be fantastical, since reality itself can be so theatrical that staging simple truths does the same work (the homespun Lincoln) --- and the honest operator who wants no part of this trade should still recognise how much of the attention economy is this floor with better upholstery.
\end{mechanism}
\begin{conditions}
Strongest where anxiety is high and verification is slow --- markets in transformation, courtship, revolutions, investment manias, grieving rooms. It weakens where feedback is fast (the fantasy meets the invoice within one cycle) and among audiences trained to price their own longing. It is at its most dangerous exactly where it is most welcome.
\end{conditions}
\begin{application}
  \item Recognise the supply: whoever makes you feel exceptional without asking anything of you is selling fantasy at retail (T-04-03).
  \item Track the specificity: a real plan has a next step with a date; a fantasy has a next feeling.
  \item When you must counter the appeal, do not argue with the dream --- offer a competing story with more romance and a working floor, or you will lose to the better fantasy every time.
  \item In your own selling, decide consciously which part of your promise is romance and which is floor, and price them separately. Customers forgive romance; they do not forgive a missing floor.
  \item Feel the pull and name the want underneath. The fantasy that catches you is a map of the thing you are not getting (T-10-02).
\end{application}
\begin{feedback}
  \item Working (as sales): the room leans in and buys the feeling before the facts. The floor under the feeling still has to exist at delivery.
  \item Working (as defence): you can state the want the pitch is selling to. Named, it loses two-thirds of its pull.
  \item Failing: the audience starts demanding the dream be enacted now. Fantasy demand ratchets; supply strains.
  \item Failing: disenchantment arrives and the room that loved you organises. Refunds in fantasy markets are paid in malice.
\end{feedback}
\begin{failure}
The oasis dries: manufacturing pleasantness is a supply business, and the operator who cannot keep topping up the romance is abandoned mid-desert --- by then the audience's tolerance for the true map has expired too. The operator also drinks their own supply; the seller of inevitable victory is the last to notice the losses. And chronic use atrophies the honesty muscle entirely: a person who never delivers disenchantment to clients eventually cannot deliver it to themselves.
\end{failure}
\begin{countermeasures}
  \item Price your longings: the fantasy that moves you fastest is pointing at the want you manage worst.
  \item Ask every seductive plan for its floor: what is true today, measurable, and unglamorous? No floor, no deal.
  \item Notice who never brings you bad news. Your information diet may be an oasis someone built around you (T-07-01).
  \item If you sell, audit your own romance-to-floor ratio quarterly. The floor is what they keep paying for; the romance is what they first pay for.
\end{countermeasures}

\sources{B3}
\seesources
\seealso{T-15-03, T-04-03, T-10-02}
